\documentclass[a4paper,10pt]{article}

\usepackage[margin=2cm]{geometry}
\usepackage{graphicx}
\usepackage{amsmath}
\usepackage{array}
\usepackage{hyperref}
\usepackage[all]{hypcap}
\usepackage{listings}
\lstdefinestyle{TerminalStyle}{
  language=bash,
  basicstyle=\small\sffamily,
  numbers=left,
  numberstyle=\tiny,
  numbersep=3pt,
  frame=tb,
  columns=fullflexible,
  linewidth=0.9\linewidth,
  xleftmargin=0.1\linewidth
}
\lstdefinestyle{HtmlStyle}{
  language=html,
  basicstyle=\small\sffamily,
  numbers=left,
  numberstyle=\tiny,
  numbersep=3pt,
  frame=tb,
  columns=fullflexible,
  linewidth=0.9\linewidth,
  xleftmargin=0.1\linewidth
}
\lstdefinestyle{OutputStyle}{
  language=html,
  basicstyle=\small\sffamily,
  frame=tb,
  columns=fullflexible,
  linewidth=0.9\linewidth,
  xleftmargin=0.1\linewidth
}

\setlength{\parindent}{0pt}
\setlength{\parskip}{1ex plus 0.5ex minus 0.2ex}

\title{\includegraphics[width=12cm]{Eeufeeslogo.jpg} \\
       Department of Computer Science \\
       University of Pretoria \\
       \vspace{0.5cm}
       Software Engineering\\
       COS301 MiniProject \\
       \vspace{0.5cm}
       \begin{large} \textbf{Team Maroon}\\ NavUP\end{large}}

\date{} 
\author{Bondjobo, Jocelyn (J) 	13232852 		\\
		Mweshi, George (G)		16394713		\\
		Letsoalo, Joseph (J)	15043844		\\
		Setoaba, Phuti (P)		13032616		\\
		Trivella, Camron (C)	14070970		\\
		Coetzer, Albert (A)		15244882		\\
}

\begin{document}
\maketitle
\thispagestyle{empty}
\clearpage

\newpage
\pagenumbering{roman}
\thispagestyle{empty}
\tableofcontents
\clearpage

\newpage
\pagenumbering{arabic}

\section{Introduction}
The client Vreda Pieterse, from the University of Pretoria has requested a system called NavUp to guide students. Students and visitors to the the campus often have to find venues to attend classes, meetings and other events or to find public services such as shops, restaurants and ablution facilities. They may want to determine the optimal route from one venue to another on a regular basis, especially at the beginning of a term when classes are presented for the first time. \\

\begin{enumerate}
\item Basic functionality
	\begin{itemize}
	\item The current location of the user should be determined both outdoors and indoors.
	\item Basic navigational functions required such searching for locations, saving locations and providing directions to a location.
	\item Front end interface
	\end{itemize}
\item Surveillance
	\begin{itemize}
	\item Functionality to provide and visualise information related to pedestrian traffic based on users last/currently using the application on campus.
	\end{itemize}
\item Targeted delivery of information
	\begin{itemize}
	\item The system can automatically load/push new information to users according to their preferences and interests.
	\end{itemize}
\item Activities
	\begin{itemize}
	\item Various activities that uses location and movement of users can be integrated into the system.
	\end{itemize}
\end{enumerate}

\section{Background}
This project was commissioned based on the following needs faced by the client:
\begin{enumerate}
\item Current situation:
	\begin{itemize}
	\item The sutents and/or visitors on campus must always ask around to people to find where a venue/building is.
	\item The sutents and/or visitors on campus had to refer to the closest map and locate their position on the map and remember direction of the building they want to go.
	\end{itemize}
\item Google map
	\begin{itemize}
	\item  Google map cannot provide navigation indoors as it is limited with the GPS satellite connection. Thus, people cannot be guided to specific venues around the campus.
	\end{itemize}
\end{enumerate}
\section{Architecture requirements}
\subsection{Access channel requirements}
\subsubsection{Human Access Channels}

\subsubsection{System Access Channels}

\subsection{Quality requirements}
\subsubsection{Performance}
\begin{enumerate}
\item \textbf{Description} \\
The aim of this application is to increase the performance of students and staff when moving to and from lecture halls and other locations on Hatfield Campus. Calculating the most efficient route between two locations while factoring in changing variables such as human traffic or events taking place in between the two locations is a key part of the performance of this application. \\
\item \textbf{Justification} \\
If the application is not performant enough then the clients using it will not be able to efficiently travel between two locations which can easily result in them getting lost or being late for lectures or meetings. \\
\item \textbf{Requirements} \\
\begin{itemize}
\item Access to WiFi routers and GPS to efficiently calculate the best route between two locations. 
\item Access to crowd sourcing to see where large clusters of students are in order to avoid traffic when calculating routes. 
\item Access to Smartphones that have the capability to use all of the features our program provides to maximise the performance of the route calculations. 
\end{itemize}
\end{enumerate}
\subsubsection{Reliability}
\begin{enumerate}
\item \textbf{Description} \\
The application needs to reliably calculate the most efficient route between two locations. The change that the program will see as more and more users begin to use it must not affect its ability to calculate the most efficient route. \\
\item \textbf{Justification} \\ 
Any application that intends to be used for a very long time needs to be able to reliably perform its given task. \\
\item \textbf{Requirements} \\
\begin{itemize}
\item Full network access of WiFi routers, especially when new ones are installed. 
\item Software managers to update the system when new buildings are built or new events are entered. 
\end{itemize}
\end{enumerate}
\subsubsection{Scalability}
\begin{enumerate}
\item \textbf{Description} \\

\item \textbf{Justification} \\
\item \textbf{Requirements}
\end{enumerate}
\subsubsection{Security}
\begin{enumerate}
\item \textbf{Description} \\

\item \textbf{Justification} \\

\item \textbf{Requirements}

\end{enumerate}
\subsubsection{Flexibility}
\begin{enumerate}
\item \textbf{Description} \\
\item \textbf{Justification} \\

\item \textbf{Requirements}
\end{enumerate}
\subsubsection{Maintainability}
\begin{enumerate}
\item \textbf{Description} \\
\item \textbf{Justification} \\
\item \textbf{Requirements}
\end{enumerate}
\subsubsection{Auditability/Monitorability}
\begin{enumerate}
\item \textbf{Description} \\
\item \textbf{Justification} \\
\item \textbf{Requirements}
\end{enumerate}
\subsubsection{Integrability}
\begin{enumerate}
\item \textbf{Description} \\
\item \textbf{Justification} \\
\item \textbf{Requirements}
\end{enumerate}
\subsubsection{Cost}
\begin{enumerate}
\item \textbf{Description} \\
\item \textbf{Justification} \\
\item \textbf{Requirements}
\end{enumerate}
\subsubsection{Usability}
\begin{enumerate}
\item \textbf{Description} \\
\item \textbf{Justification} \\
\item \textbf{Requirements}
\end{enumerate}

\subsection{Integration requirements}
The Android application will be built to work on all Android version 4.0.  and upward devices. The mobile application on the Android device should be sensitive to the amount of data that is transmitted over the network due to the high cost of mobile internet in South Africa. To allow for minimal communication but maximum functionality, the MVC based design will be followed in the Android application design. 
\subsection{Architecture constraints}
\begin{itemize}
\item Dependency Injection (DI) will be used because it eases software testing and re-usability of software by using a design based on independent classes/components.
\item Unit Testing to automate the following tasks: compilation of Java source code, running test cases and generating project documentation.
\end{itemize}
\subsubsection{Architectural Patterns}

\clearpage
\section{Open Issues}
\subsection {GitHub Repository}
Team Maroon Repository: \url{https://github.com/josephbondjobo/COS301-Maroon-Team/tree/develop}

This repository contains:
\begin{enumerate}
\item All work done by team members.
\item \href{https://github.com/josephbondjobo/COS301-Maroon-Team/blob/develop/Intro/Contributors.md}{CONTRIBUTORS.md} file outlining which members where involved in which phases of the project.
\end{enumerate}



\newpage
\clearpage
\addcontentsline{toc}{section}{References}

\end{document}
